\documentclass[11pt]{article}

\usepackage{amsmath, amssymb, amsthm}
\usepackage[margin=1in]{geometry}
\usepackage{booktabs}
\usepackage{enumitem}

\DeclareMathOperator{\disc}{disc}

\title{Discrepancy Problems: Definition and Context}
\author{}
\date{}

\begin{document}

\maketitle

\section{The General Setup}

A \textbf{discrepancy problem} is a class of problems in combinatorics
and discrepancy theory that asks:

\begin{quote}
\itshape
Given a finite set system, can we color the ground set with two colors
so that every set in the system is as ``balanced'' as possible --- and
how balanced can we guarantee it to be?
\end{quote}

More precisely, one is given:

\begin{itemize}
  \item A \textbf{ground set} \(X\) (usually finite, \(|X| = n\)),
  \item A \textbf{set system} \(\mathcal{F} \subseteq \mathcal{P}(X)\)
        --- a family of subsets of \(X\),
  \item A \textbf{two-coloring} \(\chi \colon X \to \{+1, -1\}\).
\end{itemize}

For each \(S \in \mathcal{F}\), the \textbf{discrepancy of \(S\)} with
respect to \(\chi\) is the \emph{imbalance}:
\[
\disc_\chi(S)
= \left| \sum_{x \in S} \chi(x) \right|
= \bigl| \#\{x \in S : \chi(x) = +1\}
       - \#\{x \in S : \chi(x) = -1\} \bigr|.
\]

The \textbf{discrepancy of the set system} \(\mathcal{F}\) is then
\[
\disc(\mathcal{F})
= \min_{\chi \colon X \to \{\pm 1\}}
  \; \max_{S \in \mathcal{F}} \;
  \disc_\chi(S).
\]

That is:

\begin{itemize}
  \item \textbf{Inner max}: over all sets in the system, how badly can a
        given coloring unbalance some set?
  \item \textbf{Outer min}: over all colorings, how well can we do in the
        worst case?
\end{itemize}

\hrulefill

\section{What Makes It a ``Problem''?}

The discrepancy problem for a given set system is the task of
determining (or bounding) \(\disc(\mathcal{F})\), and often of
explicitly constructing an optimal or near-optimal coloring.

There are several standard variants:

\begin{center}
\begin{tabular}{llll}
\toprule
\textbf{Variant} & \textbf{Ground set} & \textbf{Set system} & \textbf{Coloring} \\
\midrule
\textbf{Combinatorial discrepancy} & Finite set & Arbitrary family & \(\pm 1\) \\
\textbf{Spencer's problem} & \([n]\) & Arbitrary family of \(m\) sets & \(\pm 1\) \\
\textbf{Balancing sequences} & \([n]\) & Intervals / prefixes & \(\pm 1\) (or \(\{-1,0,+1\}\)) \\
\textbf{Geometric discrepancy} & Points in \(\mathbb{R}^d\) & Half-spaces, boxes, etc. & \(\pm 1\) \\
\textbf{Arithmetic discrepancy} & \(\mathbb{Z}\) or \(\mathbb{Z}_n\) & Arithmetic progressions & \(\pm 1\) \\
\bottomrule
\end{tabular}
\end{center}

\hrulefill

\section{The Key Tension}

The discrepancy problem captures a fundamental tension:

\begin{itemize}
  \item If we color too \textbf{uniformly} (e.g., half \(+1\), half
        \(-1\)), small sets might still be very unbalanced (e.g., a
        singleton is always \(\pm 1\)).
  \item If we try to \textbf{balance every small set}, we may be forced
        into a coloring that unbalances some large set badly.
\end{itemize}

The discrepancy measures the \emph{optimal worst-case tradeoff}.

\hrulefill

\section{Canonical Examples}

\textbf{1. The complete set system \(\mathcal{P}(X)\).}
As shown earlier, \(\disc(\mathcal{P}(X)) = \lceil n/2 \rceil\). The
worst sets are the monochromatic ones --- any coloring has a color class
of size at least \(\lceil n/2 \rceil\), and that class itself is a set in
the system.

\medskip
\textbf{2. Intervals on a line.}
For \(\mathcal{F} = \{\text{intervals in } [n]\}\), the discrepancy is
much smaller --- it grows like \(\log n\) (a classical result related to
the Erdős--Szekeres theorem and to the ``Steinitz lemma'' in one
dimension).

\medskip
\textbf{3. Arbitrary families of \(m\) sets (Spencer's theorem).}
For any family of \(m\) subsets of an \(n\)-element set, there exists a
coloring with discrepancy \(O(\sqrt{m \log(2n/m)})\). This is the
celebrated \textbf{``six standard deviations suffice''} theorem of
Spencer (1985).

\medskip
\textbf{4. Arithmetic progressions (Erdős discrepancy problem).}
For colorings of the integers and the system of all finite arithmetic
progressions, the discrepancy is \emph{unbounded} --- this was the
Erdős discrepancy problem, resolved affirmatively by Terence Tao in
2015.

\hrulefill

\section{Summary}

A discrepancy problem is:

\begin{quote}
\itshape
The problem of finding a two-coloring of a ground set that minimizes the
maximum imbalance over a prescribed family of subsets --- and of
determining how small that minimum can be made.
\end{quote}

It sits at the intersection of combinatorics, probability, geometry, and
number theory, and its central question is always the same: \emph{how
evenly can we distribute two colors across all sets in a given family?}

\end{document}
